\documentclass[12pt, a4paper]{article}
\usepackage[a4paper, margin=1in, headheight=24.64995pt]{geometry}
\usepackage[utf8]{inputenc}
\usepackage[T1]{fontenc}
\usepackage{lmodern}
\usepackage{graphicx}
\usepackage{titlesec}
\usepackage{fancyhdr}
\usepackage{enumitem}
\usepackage{listings}
\usepackage{amsmath}
\usepackage{tcolorbox}
\usepackage{hyperref}
\usepackage{xcolor}
\usepackage{array}
\usepackage{booktabs}
\usepackage{longtable}
\usepackage{multirow}
\usepackage{tabularx}
\usepackage{multicol}
\usepackage{listings}
\usepackage{pgfplots}
\pgfplotsset{compat=1.18}

% رنگ‌ها
\definecolor{myblue}{RGB}{0, 112, 192}
\definecolor{mygreen}{RGB}{0, 176, 80}
\definecolor{myorange}{RGB}{255, 102, 0}
\definecolor{myred}{RGB}{192, 0, 0}
\definecolor{mygray}{RGB}{240, 240, 240}

% تنظیمات صفحه
\pagestyle{fancy}
\fancyhf{}
\fancyhead[L]{\textbf{LPIC-1 Certification Notes}}
\fancyhead[R]{\textbf{Exam 101-500 \& 102-500}}
\fancyfoot[C]{\thepage}
\renewcommand{\headrulewidth}{1pt}
\renewcommand{\footrulewidth}{0.4pt}

% تنظیمات بخش‌ها
\titleformat{\section}
{\normalfont\Large\bfseries\color{myblue}}
{\thesection}{1em}{}
\titlespacing{\section}{0pt}{12pt}{6pt}

\titleformat{\subsection}
{\normalfont\large\bfseries\color{mygreen!80!black}}
{\thesubsection}{1em}{}
\titlespacing{\subsection}{0pt}{10pt}{4pt}

\titleformat{\subsubsection}
{\normalfont\normalsize\bfseries\color{myorange}}
{\thesubsubsection}{1em}{}
\titlespacing{\subsubsection}{0pt}{8pt}{2pt}

% تنظیمات لیست
\setlist[itemize]{leftmargin=*, label=$\bullet$}
\setlist[enumerate]{leftmargin=*, label=\arabic*.}

% تنظیمات کد
\lstset{
    basicstyle=\ttfamily\footnotesize,
    breaklines=true,
    frame=single,
    backgroundcolor=\color{mygray},
    tabsize=2,
    numbers=left,
    numberstyle=\tiny\color{gray},
    keywordstyle=\color{blue},
    commentstyle=\color{green!50!black},
    stringstyle=\color{orange},
    showstringspaces=false
}

% عنوان
\title{
    \textbf{\Huge LPIC-1: Linux Administrator} \\
    \vspace{0.3cm}
    \Large Complete Study Guide \& Reference Manual \\
    \vspace{0.2cm}
    \normalsize Based on Linux Professional Institute Certification
}
\author{
    \textbf{Student Name} \\
    \small \texttt{student@email.com}
}
\date{
    \textbf{\today} \\
    \small Version 2.0
}

\begin{document}

% صفحه جلد
\begin{titlepage}
    \centering
    \vspace*{2cm}
    
    % لوگو
    \begin{tikzpicture}
        \fill[myblue] (0,0) circle (1.5cm);
        \node[white, font=\Huge\bfseries] at (0,0) {L};
        \node[white, font=\Large] at (0,-0.3) {PIC-1};
    \end{tikzpicture}
    
    \vspace{1.5cm}
    
    {\Huge \bfseries Linux Professional Institute} \\
    {\Huge \bfseries Certification Level 1}
    
    \vspace{0.5cm}
    
    {\Large Complete Study Guide \& Reference Manual}
    
    \vspace{1.5cm}
    
    \begin{tcolorbox}[
        width=0.8\textwidth,
        colback=myblue!5,
        colframe=myblue,
        arc=0mm,
        boxrule=1pt,
        title=\textbf{Exam Objectives Coverage}
    ]
        \begin{multicols}{2}
            \begin{itemize}
                \item[$\checkmark$] System Architecture
                \item[$\checkmark$] Linux Installation
                \item[$\checkmark$] GNU and Unix Commands
                \item[$\checkmark$] Devices, Filesystems
                \item[$\checkmark$] Shell Scripting
                \item[$\checkmark$] User Interfaces
                \item[$\checkmark$] Administrative Tasks
                \item[$\checkmark$] Essential System Services
                \item[$\checkmark$] Networking Fundamentals
                \item[$\checkmark$] Security
            \end{itemize}
        \end{multicols}
    \end{tcolorbox}
    
    \vfill
    
    {\large \textbf{Exam Codes:} 101-500 \& 102-500} \\
    \vspace{0.2cm}
    {\large \textbf{Last Updated:} \today} \\
    \vspace{0.5cm}
    {\small \textit{For educational purposes only}}
\end{titlepage}

\tableofcontents
\newpage

% معرفی
\section{Introduction to LPIC-1 Certification}
\label{sec:intro}

\begin{tcolorbox}[
    colback=myblue!5,
    colframe=myblue,
    title=\textbf{About LPIC-1 Certification}
]
    LPIC-1 is the first certification in the Linux Professional Institute's multi-level program. It validates the ability to perform maintenance tasks on the command line, install and configure a computer running Linux, and configure basic networking.
\end{tcolorbox}

\subsection{Exam Details}
\begin{table}[h]
    \centering
    \begin{tabularx}{\textwidth}{|l|X|}
        \hline
        \textbf{Parameter} & \textbf{Details} \\
        \hline
        Exam Code & 101-500 and 102-500 \\
        \hline
        Number of Questions & 60 per exam \\
        \hline
        Duration & 90 minutes per exam \\
        \hline
        Passing Score & 500/800 (62.5\%) \\
        \hline
        Required Score & Must pass both exams \\
        \hline
        Validity & 5 years \\
        \hline
    \end{tabularx}
    \caption{LPIC-1 Exam Specifications}
\end{table}

\newpage

% بخش 1: معماری سیستم
\section{System Architecture}
\label{sec:system-arch}

\subsection{Determine and Configure Hardware Settings}

\begin{tcolorbox}[colback=mygreen!5, colframe=mygreen, title=\textbf{Key Commands}]
    \begin{lstlisting}[language=bash]
# Hardware Information
lspci          # List PCI devices
lsusb          # List USB devices
lsblk          # List block devices
lscpu          # Show CPU information
free -h        # Memory usage
dmidecode      # DMI table decoder
    \end{lstlisting}
\end{tcolorbox}

\subsubsection*{Important Files}
\begin{itemize}
    \item \texttt{/proc/cpuinfo} - CPU information
    \item \texttt{/proc/meminfo} - Memory information
    \item \texttt{/proc/interrupts} - IRQ assignments
    \item \texttt{/proc/ioports} - I/O port assignments
    \item \texttt{/proc/dma} - DMA channel assignments
\end{itemize}

\subsection{Boot the System}
\begin{table}[h]
    \centering
    \begin{tabular}{|l|l|l|}
        \hline
        \textbf{Bootloader} & \textbf{Configuration File} & \textbf{Command} \\
        \hline
        GRUB Legacy & /boot/grub/menu.lst & grub-install \\
        \hline
        GRUB 2 & /boot/grub/grub.cfg & update-grub \\
        \hline
        LILO & /etc/lilo.conf & lilo \\
        \hline
    \end{tabular}
    \caption{Linux Bootloaders}
\end{table}

\newpage

% بخش 2: نصب لینوکس و مدیریت پکیج
\section{Linux Installation and Package Management}
\label{sec:install-pkg}

\subsection{Design Hard Disk Layout}
\begin{figure}[h]
    \centering
    \begin{tikzpicture}[
        box/.style={draw, rectangle, minimum width=2cm, minimum height=1cm, align=center},
        arrow/.style={->, >=stealth, thick}
    ]
        \node[box, fill=myblue!20] (root) {/ (root)};
        \node[box, fill=mygreen!20, below left=1cm of root] (home) {/home};
        \node[box, fill=myorange!20, below=1cm of root] (var) {/var};
        \node[box, fill=myred!20, below right=1cm of root] (tmp) {/tmp};
        \node[box, fill=gray!20, above=1cm of root] (boot) {/boot};
        
        \draw[arrow] (root) -- (home);
        \draw[arrow] (root) -- (var);
        \draw[arrow] (root) -- (tmp);
        \draw[arrow] (boot) -- (root);
    \end{tikzpicture}
    \caption{Typical Linux Filesystem Hierarchy}
\end{figure}

\subsection{Install Boot Manager}
\begin{lstlisting}[language=bash, caption=GRUB 2 Installation]
# Install GRUB to MBR
grub-install /dev/sda

# Update GRUB configuration
update-grub

# For UEFI systems
grub-install --target=x86_64-efi --efi-directory=/boot/efi
\end{lstlisting}

\subsection{Manage Shared Libraries}
\begin{itemize}
    \item \texttt{ldd} - Print shared library dependencies
    \item \texttt{ldconfig} - Configure dynamic linker run-time bindings
    \item \texttt{/etc/ld.so.conf} - Configuration file for ldconfig
    \item \texttt{LD\_LIBRARY\_PATH} - Environment variable for libraries
\end{itemize}

\subsection{Package Management}
\subsubsection*{Debian/Ubuntu (APT)}
\begin{lstlisting}[language=bash]
apt update                 # Update package list
apt upgrade                # Upgrade packages
apt install package        # Install package
apt remove package         # Remove package
apt search keyword         # Search for packages
apt-cache show package     # Show package info
dpkg -i package.deb        # Install .deb file
dpkg -l                    # List installed packages
\end{lstlisting}

\subsubsection*{RHEL/CentOS/Fedora (YUM/DNF)}
\begin{lstlisting}[language=bash]
yum install package        # Install package
yum remove package         # Remove package
yum update                 # Update packages
yum search keyword         # Search packages
yum info package           # Package information
rpm -i package.rpm         # Install .rpm file
rpm -qa                    # List installed packages
rpm -ql package            # List package files
\end{lstlisting}

\newpage

% بخش 3: کامندهای GNU و Unix
\section{GNU and Unix Commands}
\label{sec:gnu-commands}

\subsection{Work on the Command Line}
\begin{table}[h]
    \centering
    \begin{tabular}{|l|l|}
        \hline
        \textbf{Command} & \textbf{Description} \\
        \hline
        \texttt{man command} & Display manual page \\
        \hline
        \texttt{info command} & Display info documentation \\
        \hline
        \texttt{whatis command} & One-line description \\
        \hline
        \texttt{which command} & Show full path of command \\
        \hline
        \texttt{whereis command} & Locate binary, source, and manual \\
        \hline
        \texttt{type command} & Display command type \\
        \hline
        \texttt{help} & Display shell built-in help \\
        \hline
    \end{tabular}
    \caption{Help and Documentation Commands}
\end{table}

\subsection{Process Text Streams Using Filters}
\begin{longtable}{|p{3cm}|p{10cm}|}
    \hline
    \textbf{Command} & \textbf{Description} \\
    \hline
    \endhead
    
    \texttt{cat} & Concatenate and display files \\
    \hline
    \texttt{tac} & Concatenate and display in reverse \\
    \hline
    \texttt{head} & Output first part of files \\
    \hline
    \texttt{tail} & Output last part of files \\
    \hline
    \texttt{cut} & Remove sections from lines \\
    \hline
    \texttt{paste} & Merge lines of files \\
    \hline
    \texttt{sort} & Sort lines of text \\
    \hline
    \texttt{uniq} & Report or omit repeated lines \\
    \hline
    \texttt{wc} & Print line, word, and byte counts \\
    \hline
    \texttt{tr} & Translate or delete characters \\
    \hline
    \texttt{expand} & Convert tabs to spaces \\
    \hline
    \texttt{unexpand} & Convert spaces to tabs \\
    \hline
    \texttt{nl} & Number lines of input \\
    \hline
    \texttt{od} & Dump files in octal and other formats \\
    \hline
    \texttt{fmt} & Simple text formatter \\
    \hline
    \texttt{pr} & Convert text files for printing \\
    \hline
\end{longtable}

\subsection{Perform Basic File Management}
\begin{lstlisting}[language=bash, caption=File Operations]
# Copy files and directories
cp file1 file2
cp -r dir1 dir2

# Move/rename files
mv oldname newname
mv file /path/to/destination/

# Remove files and directories
rm file
rm -r directory
rm -f file  # force removal

# Create and remove directories
mkdir newdir
mkdir -p path/to/newdir
rmdir emptydir

# Create links
ln -s target linkname  # symbolic link
ln target linkname     # hard link
\end{lstlisting}

\subsection{Use Streams, Pipes, and Redirects}
\begin{center}
\begin{tikzpicture}[
    node distance=1.5cm,
    box/.style={draw, rectangle, minimum width=2.5cm, minimum height=1.2cm},
    arrow/.style={->, >=stealth, thick}
]
    \node[box] (input) {Input};
    \node[box, right=of input] (process) {Process};
    \node[box, right=of process] (output) {Output};
    \node[box, below=0.8cm of process] (file) {File};
    
    \draw[arrow] (input) -- (process);
    \draw[arrow] (process) -- (output);
    \draw[arrow] (process) -- node[right] {> or >>} (file);
    \draw[arrow] (file) -- node[left] {<} (process);
\end{tikzpicture}
\end{center}

\begin{lstlisting}[language=bash, caption=I/O Redirection Examples]
# Standard output redirection
command > file        # Overwrite file
command >> file       # Append to file

# Standard error redirection
command 2> error.log  # Redirect errors
command 2>&1          # Merge stderr with stdout

# Input redirection
command < input.txt

# Pipes
command1 | command2 | command3

# Here document
cat << EOF
This is a
multiline
text
EOF
\end{lstlisting}

\subsection{Create, Monitor, and Kill Processes}
\begin{table}[h]
    \centering
    \begin{tabular}{|l|l|l|}
        \hline
        \textbf{Signal} & \textbf{Number} & \textbf{Description} \\
        \hline
        SIGHUP & 1 & Hangup (reload configuration) \\
        \hline
        SIGINT & 2 & Interrupt (Ctrl+C) \\
        \hline
        SIGKILL & 9 & Kill (cannot be caught/ignored) \\
        \hline
        SIGTERM & 15 & Terminate (graceful shutdown) \\
        \hline
        SIGSTOP & 19 & Stop (pause process) \\
        \hline
        SIGCONT & 18 & Continue (resume process) \\
        \hline
    \end{tabular}
    \caption{Common Linux Signals}
\end{table}

\subsubsection*{Process Management Commands}
\begin{lstlisting}[language=bash]
# Process viewing
ps aux          # Show all processes
ps -ef          # Full format listing
top             # Interactive process viewer
htop            # Enhanced top (if installed)
pstree          # Show process tree

# Process control
kill PID        # Send TERM signal
kill -9 PID     # Force kill
killall name    # Kill all processes with name
pkill pattern   # Kill by pattern

# Job control
jobs            # List jobs
fg %n           # Bring job n to foreground
bg %n           # Send job n to background
Ctrl+Z          # Suspend foreground job
Ctrl+C          # Interrupt process
\end{lstlisting}

\subsection{Modify Process Execution Priorities}
\begin{itemize}
    \item \texttt{nice} - Run a program with modified scheduling priority
    \item \texttt{renice} - Alter priority of running process
    \item Priority range: -20 (highest) to 19 (lowest)
    \item Only root can set negative nice values
\end{itemize}

\begin{lstlisting}[language=bash, caption=Process Priority Examples]
# Start process with low priority
nice -n 10 command

# Start process with high priority (root only)
nice -n -10 command

# Change priority of running process
renice -5 PID
\end{lstlisting}

\newpage

% بخش 4: دستگاه‌ها، فایل‌سیستم‌ها، و استاندارد FHS
\section{Devices, Linux Filesystems, Filesystem Hierarchy Standard}
\label{sec:filesystems}

\subsection{Create Partitions and Filesystems}
\begin{lstlisting}[language=bash, caption=Partition Management]
# List disks and partitions
fdisk -l
lsblk
parted -l

# Create partition with fdisk
fdisk /dev/sda
# Commands in fdisk:
# n - new partition
# p - print partition table
# d - delete partition
# w - write changes and exit
# q - quit without saving

# Create partition with parted
parted /dev/sda
mkpart primary ext4 1MiB 1GiB

# Create filesystem
mkfs.ext4 /dev/sda1
mkfs.xfs /dev/sda2
mkfs.vfat /dev/sda3

# Check filesystem
fsck /dev/sda1
e2fsck /dev/sda1
\end{lstlisting}

\subsection{Maintain the Integrity of Filesystems}
\begin{itemize}
    \item \texttt{du} - Estimate file space usage
    \item \texttt{df} - Report filesystem disk space usage
    \item \texttt{dumpe2fs} - Dump ext2/ext3/ext4 filesystem information
    \item \texttt{tune2fs} - Adjust tunable filesystem parameters
    \item \texttt{debugfs} - ext2/ext3/ext4 filesystem debugger
    \item \texttt{xfs\_info} - XFS filesystem information
\end{itemize}

\subsection{Control Mounting and Unmounting of Filesystems}
\begin{lstlisting}[language=bash, caption=Mount Operations]
# Mount filesystem
mount /dev/sda1 /mnt
mount -t ext4 /dev/sda1 /mnt
mount -o ro /dev/sda1 /mnt  # Read-only

# Unmount filesystem
umount /mnt
umount /dev/sda1

# List mounted filesystems
mount
cat /proc/mounts
findmnt

# Mount at boot (add to /etc/fstab)
/dev/sda1   /mnt/data   ext4   defaults   0   2

# Mount all from fstab
mount -a
\end{lstlisting}

\subsection{Manage Disk Quotas}
\begin{enumerate}
    \item Install quota tools: \texttt{apt-get install quota}
    \item Enable quotas in /etc/fstab: \texttt{usrquota,grpquota}
    \item Remount filesystem: \texttt{mount -o remount /}
    \item Create quota files: \texttt{quotacheck -cugm /}
    \item Turn on quotas: \texttt{quotaon /}
    \item Set quotas: \texttt{edquota username}
\end{enumerate}

\subsection{Manage File Permissions and Ownership}
\begin{table}[h]
    \centering
    \begin{tabular}{|c|l|}
        \hline
        \textbf{Permission} & \textbf{Description} \\
        \hline
        777 & rwxrwxrwx (all permissions for all) \\
        \hline
        755 & rwxr-xr-x (owner: all, group/others: read+execute) \\
        \hline
        644 & rw-r--r-- (owner: read+write, others: read) \\
        \hline
        600 & rw------- (owner: read+write, others: none) \\
        \hline
        400 & r-------- (owner: read only) \\
        \hline
        775 & rwxrwxr-x (owner/group: all, others: read+execute) \\
        \hline
    \end{tabular}
    \caption{Common Permission Numeric Codes}
\end{table}

\begin{lstlisting}[language=bash, caption=Permission Commands]
# Change permissions
chmod 755 file
chmod u+x file            # Add execute for user
chmod g-w file            # Remove write for group
chmod o=r file            # Set others to read only
chmod a+x file            # Add execute for all

# Change ownership
chown user file
chown user:group file
chown -R user directory   # Recursive

# Change group
chgrp group file
chgrp -R group directory  # Recursive

# Special permissions
chmod +s file             # Set SUID/SGID
chmod +t directory        # Sticky bit
\end{lstlisting}

\subsection{Create and Change Hard and Symbolic Links}
\begin{table}[h]
    \centering
    \begin{tabular}{|l|l|l|}
        \hline
        \textbf{Feature} & \textbf{Hard Link} & \textbf{Symbolic Link} \\
        \hline
        Created with & \texttt{ln target link} & \texttt{ln -s target link} \\
        \hline
        Inode & Same as original & Different from original \\
        \hline
        Cross-filesystem & No & Yes \\
        \hline
        Directory linking & No & Yes \\
        \hline
        After original deleted & Still works & Broken link \\
        \hline
        File size & Same as original & Path length \\
        \hline
    \end{tabular}
    \caption{Hard Links vs Symbolic Links}
\end{table}

\newpage

% بخش 5: اسکریپت‌نویسی شل و مدیریت داده‌ها
\section{Shell Scripting and Data Management}
\label{sec:shell-scripting}

\subsection{Customize and Use the Shell Environment}
\begin{lstlisting}[language=bash, caption=Shell Configuration]
# Environment variables
echo $PATH              # Command search path
echo $HOME              # Home directory
echo $USER              # Current user
echo $SHELL             # Current shell
echo $PS1               # Prompt string
echo $PWD               # Present working directory

# Set environment variables
export VAR=value        # For current shell and children
VAR=value               # For current shell only

# Shell configuration files
~/.bashrc               # Non-login interactive shells
~/.bash_profile         # Login shells
~/.profile              # Alternative profile
/etc/bashrc             # System-wide bashrc
/etc/profile            # System-wide profile

# Aliases
alias ll='ls -la'
alias rm='rm -i'
unalias ll              # Remove alias
\end{lstlisting}

\subsection{Customize or Write Simple Scripts}
\begin{lstlisting}[language=bash, caption=Basic Shell Script Template]
#!/bin/bash
# Script: backup.sh
# Description: Backup script example
# Author: Your Name
# Date: $(date)

# Configuration
BACKUP_DIR="/backup"
SOURCE_DIR="/home/user/documents"
DATE=$(date +%Y%m%d_%H%M%S)
LOG_FILE="/var/log/backup.log"

# Functions
log_message() {
    echo "[$(date)] $1" >> "$LOG_FILE"
    echo "$1"
}

# Check if running as root
if [[ $EUID -ne 0 ]]; then
    echo "This script must be run as root" >&2
    exit 1
fi

# Check if source directory exists
if [ ! -d "$SOURCE_DIR" ]; then
    log_message "ERROR: Source directory $SOURCE_DIR not found"
    exit 1
fi

# Create backup directory if it doesn't exist
mkdir -p "$BACKUP_DIR"

# Perform backup
log_message "Starting backup of $SOURCE_DIR"
tar -czf "$BACKUP_DIR/backup_$DATE.tar.gz" "$SOURCE_DIR"

# Check if backup was successful
if [ $? -eq 0 ]; then
    log_message "Backup completed successfully: backup_$DATE.tar.gz"
    # Remove backups older than 30 days
    find "$BACKUP_DIR" -name "*.tar.gz" -mtime +30 -delete
else
    log_message "ERROR: Backup failed"
    exit 1
fi

log_message "Backup process completed"
\end{lstlisting}

\newpage

% بخش 6: کاربران و گروه‌ها
\section{User and Group Administration}
\label{sec:user-admin}

\subsection{Manage User and Group Accounts}
\begin{lstlisting}[language=bash, caption=User Management Commands]
# Add user
useradd username
useradd -m username      # Create home directory
useradd -g groupname username  # Primary group
useradd -G group1,group2 username  # Supplementary groups
useradd -s /bin/bash username  # Specify shell

# Modify user
usermod -aG groupname username  # Add to group
usermod -s /bin/sh username     # Change shell
usermod -L username            # Lock account
usermod -U username            # Unlock account

# Delete user
userdel username
userdel -r username     # Remove home directory and mail

# Manage groups
groupadd groupname
groupdel groupname
gpasswd -a username groupname  # Add user to group
gpasswd -d username groupname  # Remove user from group

# Password management
passwd username
passwd -l username      # Lock password
passwd -u username      # Unlock password
chage username          # Change password aging
\end{lstlisting}

\subsection{Configure User Environment}
\begin{itemize}
    \item \texttt{/etc/passwd} - User account information
    \item \texttt{/etc/shadow} - Secure user password information
    \item \texttt{/etc/group} - Group information
    \item \texttt{/etc/gshadow} - Secure group password information
    \item \texttt{/etc/skel/} - Skeleton directory for new users
    \item \texttt{/etc/login.defs} - Configuration for useradd
\end{itemize}

\newpage

% بخش 7: امنیت
\section{Security}
\label{sec:security}

\subsection{Perform Security Administration Tasks}
\begin{itemize}
    \item \texttt{su} - Switch user
    \item \texttt{sudo} - Execute commands as another user
    \item \texttt{visudo} - Edit sudoers file safely
    \item \texttt{chroot} - Run command with special root directory
    \item \texttt{setuid/setgid} - Special file permissions
    \item \texttt{umask} - Default file permissions
\end{itemize}

\subsection{Set Up Host Security}
\begin{lstlisting}[language=bash, caption=Security Commands]
# SSH configuration
/etc/ssh/sshd_config    # SSH server configuration
ssh-keygen              # Generate SSH keys
ssh-copy-id             # Copy SSH key to remote host

# Firewall management (iptables)
iptables -L             # List rules
iptables -A INPUT -p tcp --dport 22 -j ACCEPT
iptables -A INPUT -j DROP

# Firewall management (firewalld)
firewall-cmd --list-all
firewall-cmd --add-port=80/tcp --permanent
firewall-cmd --reload

# SELinux management
getenforce              # Check SELinux status
sestatus                # Detailed SELinux status
setenforce 0            # Disable SELinux (temporary)
setenforce 1            # Enable SELinux (temporary)
\end{lstlisting}

\newpage

% بخش 8: شبکه
\section{Networking}
\label{sec:networking}

\subsection{Fundamentals of Internet Protocols}
\begin{itemize}
    \item TCP/IP Protocol Suite
    \item IPv4 and IPv6 Addressing
    \item Subnetting and CIDR Notation
    \item Common Ports and Services
    \begin{itemize}
        \item 22 - SSH
        \item 80 - HTTP
        \item 443 - HTTPS
        \item 53 - DNS
        \item 25 - SMTP
        \item 110 - POP3
        \item 143 - IMAP
    \end{itemize}
\end{itemize}

\subsection{Basic Network Configuration}
\begin{lstlisting}[language=bash, caption=Network Commands]
# Network interface configuration
ifconfig eth0 192.168.1.100 netmask 255.255.255.0
ip addr add 192.168.1.100/24 dev eth0

# Network interface control
ifconfig eth0 down
ifconfig eth0 up
ip link set eth0 down
ip link set eth0 up

# Routing configuration
route add default gw 192.168.1.1
ip route add default via 192.168.1.1
route -n
ip route show

# DNS configuration
/etc/resolv.conf        # DNS resolver configuration
/etc/hosts              # Local hostname resolution
\end{lstlisting}

\subsection{Troubleshooting Network Issues}
\begin{lstlisting}[language=bash]
ping google.com          # Test connectivity
traceroute google.com    # Trace route
netstat -tulpn           # Show open ports
ss -tulpn                # Alternative to netstat
nslookup google.com      # DNS lookup
dig google.com           # DNS query
host google.com          # Hostname lookup
\end{lstlisting}

\newpage

% بخش 9: سرویس‌های ضروری
\section{Essential System Services}
\label{sec:services}

\subsection{System Logging}
\begin{itemize}
    \item \texttt{rsyslog} - System logging daemon
    \item \texttt{/etc/rsyslog.conf} - Configuration file
    \item \texttt{/var/log/} - Log directory
    \begin{itemize}
        \item \texttt{messages} - General system messages
        \item \texttt{auth.log} - Authentication logs
        \item \texttt{syslog} - System logs
        \item \texttt{kern.log} - Kernel logs
        \item \texttt{dpkg.log} - Package manager logs
    \end{itemize}
    \item \texttt{logrotate} - Rotate, compress, and mail logs
\end{itemize}

\subsection{Time Synchronization}
\begin{itemize}
    \item \texttt{ntpd} - Network Time Protocol daemon
    \item \texttt{chronyd} - Chrony NTP daemon
    \item \texttt{date} - Display or set system date and time
    \item \texttt{timedatectl} - Control system time and date
    \item \texttt{hwclock} - Hardware clock access
\end{itemize}

\subsection{Mail Transfer Agent (MTA) Basics}
\begin{itemize}
    \item \texttt{postfix} - Mail transfer agent
    \item \texttt{sendmail} - Alternative MTA
    \item \texttt{mail} - Send and receive mail
    \item \texttt{mailq} - Print mail queue
\end{itemize}

\newpage

% بخش پایانی
\section{Appendix: Exam Tips and Resources}
\label{sec:appendix}

\begin{tcolorbox}[
    colback=myorange!5,
    colframe=myorange,
    title=\textbf{Exam Preparation Tips}
]
    \begin{itemize}
        \item \textbf{Practice on real Linux systems} - Use VirtualBox or VMware
        \item \textbf{Learn command line options} - Focus on commonly used switches
        \item \textbf{Understand configuration files} - Know location and syntax
        \item \textbf{Time management} - Each question has about 1.5 minutes
        \item \textbf{Read questions carefully} - Pay attention to details
        \item \textbf{Use elimination method} - Remove obviously wrong answers
    \end{itemize}
\end{tcolorbox}

\subsection{Recommended Resources}
\begin{itemize}
    \item \textbf{Official LPIC-1 Objectives}: https://www.lpi.org/our-certifications/exam-101-objectives
    \item \textbf{LPIC-1 Study Guide} by Christine Bresnahan and Richard Blum
    \item \textbf{Linux Command Line and Shell Scripting Bible} by Richard Blum
    \item \textbf{Online Practice Exams}: Boson, ExamCompass, ExamTopics
    \item \textbf{Linux Documentation Project}: https://tldp.org/
\end{itemize}

\subsection{Command Quick Reference}
\begin{table}[h]
    \centering
    \begin{tabular}{|l|l|}
        \hline
        \textbf{Category} & \textbf{Key Commands} \\
        \hline
        File Management & ls, cp, mv, rm, find, chmod, chown \\
        \hline
        Process Management & ps, top, kill, nice, jobs \\
        \hline
        Text Processing & grep, sed, awk, cut, sort, uniq \\
        \hline
        Networking & ping, ifconfig, netstat, ssh, scp \\
        \hline
        Package Management & apt, yum, dpkg, rpm \\
        \hline
        System Info & uname, df, du, free, lspci, lsusb \\
        \hline
        User Management & useradd, usermod, passwd, groupadd \\
        \hline
        Disk Management & fdisk, mkfs, mount, fsck \\
        \hline
    \end{tabular}
    \caption{LPIC-1 Essential Commands}
\end{table}

\begin{center}
\vspace{1cm}
\hrule
\vspace{0.5cm}
\textbf{\Large Good Luck with Your LPIC-1 Certification!} \\
\vspace{0.2cm}
\textit{Remember: Practice is key to success in Linux administration.}
\vspace{0.5cm}
\hrule
\end{center}

\end{document}